\scanpage{321}\mspage{45}
% archive sheet number: 320

(On désignera la condition obtenue sur le couple $W\supset W_0$ de loc fond. par \deleted{la condition} $D_{W_0}$ \deleted{$\unsure{S|D}$}).

Ainsi on a (comme $\mathrm{Parf}_{W_{\unsure{0|\bullet}}}\subset\mathrm{Parf}_{W}$)
\[
 D=D_W\Longrightarrow D_{W_0}, \tag*{(4.4.4)}
\]
et si $W'_0\subset W_0$, on aura \deleted{$D_{W_0}\Longrightarrow D_{W'_0}$}.
\[
 D_{W_0}\Longrightarrow D_{W'_0}. \tag*{(4.4.5)}
\]

\uline{Corollaire} 4.4.5. Désignons par $A,B,C$ respectivement les conditions $\mathrm{Loc}(7\unsure{a|\bullet}),\mathrm{Loc}(7b),\mathrm{Loc}(7\unsure{c|b})$. On a alors les implications
\[
 A\Longrightarrow B\Longrightarrow C. \tag*{(4.4.\deleted{\unsure{8}}\,\unsure{5.1|8.1})}
\]
De plus, \deleted{\unsure{si}} on désigne (pour un ens. $\unsure{\mathcal{F}|F}$ de flèches de $\mathrm{Cat}$) par $\mathrm{Comp}(\mathcal{F})$ la condition sur $\mathcal{F}$ que $\mathcal{F}$ soit stable par composition, on a les équivalences
\[
 A\Longleftrightarrow(C+\mathrm{Comp}(\mathcal{F}_{W}))
 \Longleftrightarrow(C+\mathrm{Comp}(\mathrm{Fibf}_{W})). \tag*{(4.4.\deleted{\unsure{1}}\,\unsure{5.2|5,2})}
\]

Dém. La première assertion est contenue dans le corollaire \unsure{4.4.7|4.4.2}, appliqué au cas $\mathcal{F}=\mathcal{F}_{W}$. D'autre part, si $\mathcal{F}_{W}$ satisfait à $\mathrm{Comp}(\mathcal{F}_{W})$, il satisfait aux conditions 1) \unsure{2)|?)} 3) du lemme de Brown 4.4.3, donc joint à $C$, il implique $A$. \deleted{\unsure{Pour}} Il reste, pour la première équivalence, à prouver que
\[
 A\Longrightarrow\mathrm{Comp}(\mathcal{F}_{W}).
\]
